\documentclass{article}

% if you need to pass options to natbib, use, e.g.:
%     \PassOptionsToPackage{numbers, compress}{natbib}
% before loading neurips_2024


% ready for submission
\usepackage{neurips_2024}


% to compile a preprint version, e.g., for submission to arXiv, add add the
% [preprint] option:
%     \usepackage[preprint]{neurips_2024}


% to compile a camera-ready version, add the [final] option, e.g.:
%     \usepackage[final]{neurips_2024}


% to avoid loading the natbib package, add option nonatbib:
%    \usepackage[nonatbib]{neurips_2024}


\usepackage[utf8]{inputenc} % allow utf-8 input
\usepackage[T1]{fontenc}    % use 8-bit T1 fonts
\usepackage{hyperref}       % hyperlinks
\usepackage{url}            % simple URL typesetting
\usepackage{graphicx}       % \resizebox for wide tables
\usepackage{booktabs}       % professional-quality tables
\usepackage{amsfonts}       % blackboard math symbols
\usepackage{nicefrac}       % compact symbols for 1/2, etc.
\usepackage{microtype}      % microtypography
\usepackage{xcolor}         % colors


\title{Zero-Shot Adaptation of Vision-Language-Action Models via World Models}


% The \author macro works with any number of authors. There are two commands
% used to separate the names and addresses of multiple authors: \And and \AND.
%
% Using \And between authors leaves it to LaTeX to determine where to break the
% lines. Using \AND forces a line break at that point. So, if LaTeX puts 3 of 4
% authors names on the first line, and the last on the second line, try using
% \AND instead of \And before the third author name.


\author{%
  David S.~Hippocampus\thanks{Use footnote for providing further information
    about author (webpage, alternative address) - \emph{not} for acknowledging
    funding agencies.} \\
  Department of Computer Science\\
  Cranberry-Lemon University\\
  Pittsburgh, PA 15213 \\
  \texttt{hippo@cs.cranberry-lemon.edu} \\
  % examples of more authors
  % \And
  % Coauthor \\
  % Affiliation \\
  % Address \\
  % \texttt{email} \\
  % \AND
  % Coauthor \\
  % Affiliation \\
  % Address \\
  % \texttt{email} \\
  % \And
  % Coauthor \\
  % Affiliation \\
  % Address \\
  % \texttt{email} \\
  % \And
  % Coauthor \\
  % Affiliation \\
  % Address \\
  % \texttt{email} \\
}


\begin{document}


\maketitle


\begin{abstract}
Vision-language-action (VLA) policies show great promise of generalisation due to their extensive training data priors, where one checkpoint can cover many instructions and different observations, but they still misbehave when the visuals or dynamics drift from training. We are exploring test-time adaptations, rather than fine-tuning, to make these VLAs even more generalisable. We start with \emph{segment-level} reranking: at each segment, sample $K$ short action chunks from a frozen SmolVLA policy, score each chunk with a JEPA-style world model by latent distance to a goal, execute the best chunk in the real simulator, and optionally apply a tiny GRPO-style preference update on top of the WM scores. The present write-up describes an early harness (logging, comparison strips, oracle-aligned goals) rather than a finished benchmark story. \textbf{TODO:} aggregate success rates, ablations, and a cleaner theoretical positioning once the pilot campaigns stabilize.
\end{abstract}



\section{Intro}
All experiments below are \textbf{Meta-World (MT1)} in simulation: same ecosystem, different \texttt{MT1(task)} strings (e.g.\ \texttt{push-v3} for the current segment-GRPO driver; other tasks use the same artifact layout where scripts exist). We compare (i) scripted oracle rollouts, (ii) frozen SmolVLA open-loop baselines, and (iii) WM-scored segment planning; no LIBERO/RoboCasa stack in this draft.

\begin{itemize}
  \item \textbf{Problem.} A frozen VLA can be wrong for a full horizon even when short chunks are locally plausible. We ask whether a \emph{test-time} reranker - a JEPA-style world model (WM) scoring candidate action chunks in latent space - improves goal progress without updating SmolVLA weights, and whether a small GRPO-style update on auxiliary parameters helps beyond pure argmax selection.

  \item \textbf{World model (JEPA-WMS).} We load a pretrained JEPA world-model bundle (Facebook \texttt{jepa-wms} codebase + checkpoint, e.g.\ \texttt{jepa\_wm\_metaworld.pth.tar}) and use it in two roles: (a) \emph{scoring} - unroll each candidate chunk in WM latent space and measure $L_2$ distance from the predicted latent prefix to a \emph{goal latent} (encode oracle RGB, optionally with proprio from the live state; configurable slice \texttt{visual} / \texttt{proprio} / \texttt{concat}); (b) \emph{diagnostics} - decode the latent trace for the selected chunk to RGB for side-by-side comparison strips vs.\ simulator frames. Rollout can be \texttt{iterative} (one WM step per env action) or \texttt{batched}; megastep stride affects how many strip columns map to how many env steps.

  \item \textbf{Phase 6  -  oracle baseline (\texttt{artifacts/phase06\_oracle\_baseline}).} Meta-World \textbf{scripted} policies (\texttt{metaworld.policies.ENV\_POLICY\_MAP}), not SmolVLA. Pipeline entrypoints: \texttt{scripts/oracle/run\_metaworld\_oracle\_eval.py} via \texttt{run\_oracle\_baseline\_eval.sh} / \texttt{pushv3\_oracle\_data\_pipeline.sh}. Outputs per run include \texttt{eval\_info.json}, \texttt{run\_manifest.json}, per-step \texttt{actions.jsonl}, and (when enabled) per-frame PNGs under \texttt{frames/episode\_XXXX/frame\_*.png}. Top-$K$ episodes are summarized to \texttt{optimal\_report.json} for downstream targeting. Segment-GRPO resolves the latest compatible run under this root (task tag in dirname) unless \texttt{--oracle-run-root} overrides.

  \item \textbf{Phase 7  -  SmolVLA baseline (\texttt{artifacts/phase07\_smolvla\_baseline}).} Frozen SmolVLA rollouts on the same Meta-World tasks via \texttt{scripts/smolvla/run\_metaworld\_smolvla\_eval.py} (LeRobot Meta-World backend; checkpoint such as \texttt{jadechoghari/smolvla\_metaworld}). Standard artifacts: \texttt{eval\_info.json}, \texttt{run\_manifest.json}, videos / action logs as configured. Top-15 oracle rows drive matched-episode campaigns under \texttt{campaigns/} so SmolVLA episodes align with phase-6 episode indices and reset seeds for apples-to-apples comparisons.

  \item \textbf{Phase 8  -  segment WM + GRPO pilot (\texttt{artifacts/phase08\_segment\_grpo\_baseline}).} \texttt{scripts/run\_segment\_grpo.py} $\rightarrow$ \texttt{rollout\_with\_chunks}: each segment samples $K$ length-$A$ action matrices from SmolVLA, scores all candidates with the WM vs.\ encoded oracle goal, executes the highest-scoring chunk in Meta-World sim (with camera/proprio parity flags aligned to WM training), logs JSON (\texttt{EpisodeLog}), and writes comparison PNGs (optional per-column overlay + footer with per-candidate integer $L_2$ traces). Optional \texttt{--train-steps} applies a tiny MLP adapter step using softmax weights from WM scores (GRPO-style group update); default pilot runs often leave this at zero.

  \item \textbf{Oracle frames in phase 8.} Goal (and optional start) pixels are \textbf{read from phase-6 PNGs}, not re-rendered: \texttt{load\_oracle\_reference\_frames} maps 1-based \texttt{--goal-frame-index} to \texttt{frame\_\{index-1\}.png} under the chosen oracle \texttt{run\_*} and episode directory; used for WM encode and reset sanity checks.

  \item \textbf{Open items (paper + experiments).}
  \begin{itemize}
    \item \textbf{TODO:} formalize segment-GRPO objective vs.\ standard GRPO and what the adapter is allowed to see.
    \item \textbf{TODO:} full quantitative table across tasks (not only \texttt{push-v3}), $K$, chunk length, WM rollout mode, and scoring latent ablation.
    \item \textbf{TODO:} strip NeurIPS instruction boilerplate below when tightening page count.
  \end{itemize}
\end{itemize}


\subsection{Evaluation and results}
\begin{itemize}
  \item Current phase-8 runs prioritize \textbf{traceability}: stitched comparison strips, per-segment JSON (\texttt{candidate\_wm\_goal\_l2\_int}, decode/scoring status fields), and matched seeds vs.\ phase-7 / phase-6 references.
  \item \textbf{TODO:} freeze primary metrics (task success, WM latent margin, human-readable strip QC rubric) and report confidence intervals over episode sets.
  \item \textbf{TODO:} one main-figure result; relegate strip galleries to appendix or supplementary.
\end{itemize}

\paragraph{Pilot (push-v3): per-megastep WM goal $L_2$ on the selected chunk.}
Table~\ref{tab:segment-grpo-wm-l2-pilot} lists the top-15 oracle-aligned Meta-World episodes (phase-7 \texttt{targets.json} ranks). For each episode we run one segment of segment-GRPO with $K{=}5$, chunk length 30, max steps 30, oracle goal frame index 25, WM stride five env steps per WM megastep. Columns \texttt{0--5}, \ldots, \texttt{25--30} are integer-rounded $L_2$ distances to the encoded oracle goal latent after each WM unroll step, logged as \texttt{d\_goal\_l2\_wm\_int} for the \emph{selected} candidate only. The \textbf{closest} column names the half-open env-step band where that distance is smallest (ties broken by the earliest band).

\begin{table}[t]
  \caption{WM integer goal $L_2$ by megastep (stride 5) for the selected chunk, push-v3 top-15 episodes.}
  \label{tab:segment-grpo-wm-l2-pilot}
  \centering
  \resizebox{\linewidth}{!}{%
  \begin{tabular}{@{}rrrcccccc@{}}
    \toprule
    Rank & Ep. & Seed & Closest & $0{:}5$ & $5{:}10$ & $10{:}15$ & $15{:}20$ & $20{:}25$ & $25{:}30$ \\
    \midrule
     1 &   3 & 1003 & $25{:}30$ & 361 & 353 & 346 & 329 & 294 & 289 \\
     2 &  53 & 1053 & $25{:}30$ & 361 & 353 & 346 & 329 & 294 & 289 \\
     3 &  16 & 1016 & $25{:}30$ & 369 & 368 & 370 & 363 & 345 & 311 \\
     4 &  44 & 1044 & $0{:}5$   & 343 & 349 & 377 & 401 & 418 & 423 \\
     5 &  39 & 1039 & $25{:}30$ & 348 & 346 & 342 & 331 & 316 & 304 \\
     6 &   8 & 1008 & $25{:}30$ & 380 & 374 & 366 & 342 & 314 & 271 \\
     7 &  58 & 1058 & $25{:}30$ & 380 & 374 & 366 & 342 & 314 & 271 \\
     8 &  10 & 1010 & $15{:}20$ & 340 & 332 & 320 & 285 & 302 & 309 \\
     9 &  42 & 1042 & $0{:}5$   & 353 & 365 & 387 & 410 & 424 & 426 \\
    10 &   9 & 1009 & $20{:}25$ & 335 & 333 & 335 & 315 & 307 & 309 \\
    11 &  59 & 1059 & $20{:}25$ & 335 & 333 & 335 & 315 & 307 & 309 \\
    12 &   5 & 1005 & $15{:}20$ & 290 & 278 & 260 & 259 & 298 & 328 \\
    13 &  55 & 1055 & $15{:}20$ & 290 & 278 & 260 & 259 & 298 & 328 \\
    14 &  33 & 1033 & $5{:}10$  & 327 & 326 & 334 & 351 & 368 & 386 \\
    15 &  28 & 1028 & $15{:}20$ & 319 & 320 & 315 & 306 & 344 & 385 \\
    \bottomrule
  \end{tabular}%
  }
\end{table}


% every seed == different starting environment of same task push-v3
% 10 different environments
% MT10
% 10 seeds, 50 episodes, avg latent distance

% 

% frameskip = 5
% 20D action into WM
% check if smovla no frameskip







\begin{itemize}
  \item \texttt{phase06\_oracle\_baseline/run\_20260411T131839Z\_ep60\_voracle\_tpush\_v3\_s1000\_r402093}: 
    from queue at \texttt{13:18:39 UTC} to finish at \texttt{13:32:26.260964 UTC} (13m 47s).
  \item \texttt{phase08\_segment\_grpo\_baseline/run\_20260416T230509Z\_all60\_f50\_k3\_s1000}: 
    from queue at \texttt{23:05:09 UTC} to finish at \texttt{00:49:23.357243 UTC} on 2026-04-17 (1h 44m 14s).
  \item \texttt{phase09\_oracle\_vs\_wm\_baseline/run\_20260416T223708Z\_ep60\_vphase09\_oracle\_vs\_wm\_tpush\_v3\_s1000\_r709580}: 
    from queue at \texttt{22:37:08 UTC} to finish at \texttt{22:45:41 UTC} (8m 33s).
\end{itemize}






\par\smallskip
\nopagebreak[4]
\begin{minipage}{\linewidth}
\raggedright

% --- INTRO ---
\noindent\textbf{Setup (policy run).} Meta-World \texttt{push-v3}; WM---goal distance is Euclidean L2 in \texttt{visual} scoring latent, integer-rounded per WM megastep (artifact \texttt{d\_goal\_l2\_wm\_int}).

Run directory: \texttt{/vol/bitbucket/aa6622/project/artifacts/phase08\_segment\_grpo\_baseline/run\_20260416T230509Z\_all60\_f50\_k3\_s1000}; episodes aggregated: 60; candidate mode: \texttt{selected}.

\noindent\textbf{Setup (oracle replay run).} Same WM and goal latent definition; actions come from recorded oracle \texttt{actions.jsonl} (phase09-style replay).

Compare directory: \texttt{/vol/bitbucket/aa6622/project/artifacts/phase09\_oracle\_vs\_wm\_baseline/run\_20260416T223708Z\_ep60\_vphase09\_oracle\_vs\_wm\_tpush\_v3\_s1000\_r709580}; paired episodes: 60.

% --- PAIRED TABLE ---
\begin{center}
\begin{tabular}{lrrrrrr}
\hline
Range & $\mu_{\mathrm{pol}}$ & SE$_{\mathrm{pol}}$ & $n_{\mathrm{pol}}$ & $\mu_{\mathrm{orc}}$ & SE$_{\mathrm{orc}}$ & $n_{\mathrm{orc}}$ \\
\hline
0:5 &  412.57 &  2.94 & 60 &  416.18 &  2.69 & 60 \\
5:10 &  407.27 &  2.97 & 60 &  414.90 &  2.60 & 60 \\
10:15 &  404.32 &  3.25 & 60 &  415.37 &  2.92 & 60 \\
15:20 &  394.22 &  3.93 & 60 &  398.47 &  3.96 & 60 \\
20:25 &  384.82 &  4.47 & 60 &  376.43 &  5.50 & 60 \\
25:30 &  372.50 &  5.17 & 60 &  360.32 &  6.92 & 60 \\
30:35 &  361.42 &  5.93 & 60 &  342.60 &  7.44 & 60 \\
35:40 &  348.38 &  5.72 & 60 &  329.67 &  7.23 & 60 \\
40:45 &  321.77 &  4.95 & 60 &  320.22 &  5.42 & 60 \\
45:50 &  294.28 &  4.89 & 60 &  329.52 &  5.17 & 60 \\
\hline
\end{tabular}
\end{center}

% --- NOTE ---
\medskip
\noindent\footnotesize\textit{Note.} Table shows cross-episode mean $\pm$ SE of per-bin values; bins are half-open env action step ranges aligned to \texttt{comparison\_wm\_env\_steps\_per\_wm\_step}. Not pixel MSE to goal image.
\noindent\footnotesize\textit{Oracle column:} WM conditioned on expert action sequence, not a second world model.
\noindent\footnotesize\textit{Policy column:} SmolVLA-selected chunk ($K{=}3$ run) per artifact defaults.

\end{minipage}
\par\smallskip
\nopagebreak[4]










% MT10 WM-goal latent L2 task-wise table: VLA and oracle delta per range.
% Requires \texttt{graphicx} for \texttt{\textbackslash resizebox}.
% Paste fragment; add booktabs to preamble if desired.

\noindent\textbf{Experiment setup.}
\begin{itemize}
\item Phase08 run: \textbf{SmolVLA policy} rollouts.
\item Phase09 run: \textbf{oracle-action replay} baseline.
\item Metric: world-model latent distance $d\_goal\_l2\_wm\_int$.
\item Binning: env-action bins 0:5, 5:10, \dots, 45:50 (10 bins).
\item For each task, rows are shown as \textbf{VLA} and \textbf{$\Delta_{\mathrm{oracle}}$} (oracle minus VLA).
\item Task names have the \texttt{-v3} suffix removed for readability.
\end{itemize}

\begin{center}
\resizebox{\textwidth}{!}{%
\begin{tabular}{llrrrrrrrrrr}
\hline
\textbf{Task} & \textbf{Run} & \textbf{0:5} & \textbf{5:10} & \textbf{10:15} & \textbf{15:20} & \textbf{20:25} & \textbf{25:30} & \textbf{30:35} & \textbf{35:40} & \textbf{40:45} & \textbf{45:50} \\
\hline
button-press-topdown & VLA & 570.92 & 554.47 & 524.98 & 495.37 & 458.78 & 415.63 & 381.60 & 350.32 & 330.50 & 361.93 \\
button-press-topdown & $\Delta_{\mathrm{oracle}}$ & -0.20 & -0.12 & +1.58 & +3.43 & +1.43 & +0.28 & -2.97 & -4.25 & +1.22 & +0.60 \\
\hline
door-open & VLA & 473.13 & 458.22 & 453.07 & 452.80 & 450.83 & 443.75 & 431.15 & 424.08 & 437.78 & 474.72 \\
door-open & $\Delta_{\mathrm{oracle}}$ & +0.00 & -0.80 & +1.18 & +1.80 & -0.45 & -2.98 & -10.23 & -20.53 & -20.82 & -10.25 \\
\hline
drawer-close & VLA & 410.33 & 407.42 & 408.95 & 416.27 & 411.33 & 394.62 & 388.43 & 387.20 & 386.02 & 390.62 \\
drawer-close & $\Delta_{\mathrm{oracle}}$ & +1.02 & +1.35 & +1.88 & +2.15 & +2.18 & +2.05 & -1.10 & +2.93 & +7.28 & +2.75 \\
\hline
drawer-open & VLA & 389.43 & 377.90 & 376.12 & 376.87 & 372.62 & 363.68 & 353.70 & 348.33 & 347.32 & 350.30 \\
drawer-open & $\Delta_{\mathrm{oracle}}$ & -1.28 & -1.93 & +0.85 & +6.57 & +15.58 & +13.70 & +7.45 & +1.37 & -3.02 & -5.13 \\
\hline
peg-insert-side & VLA & 465.98 & 465.92 & 467.95 & 463.37 & 456.08 & 445.95 & 434.20 & 423.33 & 416.35 & 412.68 \\
peg-insert-side & $\Delta_{\mathrm{oracle}}$ & -1.73 & -1.58 & +0.37 & +3.45 & +1.28 & -0.37 & -2.30 & -7.97 & -20.30 & -27.35 \\
\hline
pick-place & VLA & 438.63 & 432.53 & 431.27 & 430.63 & 433.02 & 434.77 & 425.22 & 397.90 & 365.05 & 346.85 \\
pick-place & $\Delta_{\mathrm{oracle}}$ & -2.08 & -3.27 & -6.45 & -6.08 & -5.43 & -5.48 & -3.37 & +0.40 & +2.58 & -1.17 \\
\hline
push & VLA & 415.18 & 412.08 & 411.30 & 404.22 & 394.48 & 380.77 & 366.05 & 351.73 & 330.98 & 312.40 \\
push & $\Delta_{\mathrm{oracle}}$ & +0.52 & +2.25 & +6.10 & +9.75 & +9.07 & +12.88 & +21.62 & +28.40 & +24.70 & +18.97 \\
\hline
reach & VLA & 489.28 & 463.28 & 447.03 & 430.03 & 410.25 & 395.38 & 385.48 & 379.00 & 378.02 & 381.02 \\
reach & $\Delta_{\mathrm{oracle}}$ & -0.60 & -2.42 & -7.08 & -9.58 & -10.80 & -10.47 & -10.25 & -10.15 & -9.77 & -8.57 \\
\hline
window-close & VLA & 576.15 & 554.33 & 523.92 & 494.13 & 471.65 & 463.00 & 467.68 & 469.35 & 471.90 & 472.45 \\
window-close & $\Delta_{\mathrm{oracle}}$ & +0.10 & +0.15 & +0.48 & +2.33 & +9.90 & +7.12 & +6.37 & +5.23 & +1.67 & +2.53 \\
\hline
window-open & VLA & 482.92 & 454.50 & 414.17 & 403.37 & 405.10 & 413.02 & 415.30 & 413.80 & 410.72 & 408.42 \\
window-open & $\Delta_{\mathrm{oracle}}$ & +0.37 & -1.58 & -2.05 & -2.82 & -0.10 & -3.35 & -6.48 & -7.00 & -5.52 & -4.25 \\
\hline
All tasks & VLA & 471.20 & 458.06 & 445.88 & 436.71 & 426.41 & 415.06 & 404.88 & 394.50 & 387.46 & 391.14 \\
All tasks & $\Delta_{\mathrm{oracle}}$ & -0.39 & -0.79 & -0.31 & +1.10 & +2.27 & +1.34 & -0.13 & -1.16 & -2.20 & -3.19 \\
\hline
\hline
\end{tabular}
}
\end{center}
\par\vspace{0.3em}
\noindent\small\textbf{How to read this table.} 7/10 are closer to oracle than policy, which is what we want. Cell values are cross-episode mean latent distances per bin. Each task has a VLA row and a \(\Delta_{\mathrm{oracle}}\) row, where \(\Delta_{\mathrm{oracle}} = \mathrm{oracle} - \mathrm{VLA}\). A positive value means oracle is farther from the world-model goal than VLA, and a negative value means oracle is closer. In this representation, larger raw values mean larger latent distance, so values closer to zero are preferable when the distance is interpreted as reconstruction error.




% MT10 WM-goal latent L2 task-wise table: VLA and oracle delta per range.
% Requires \texttt{graphicx} for \texttt{\textbackslash resizebox}.
% Paste fragment; add booktabs to preamble if desired.

\noindent\textbf{Experiment setup.}
\begin{itemize}
\setlength{\itemsep}{0.25em}
\setlength{\parsep}{0pt}
\setlength{\topsep}{0.25em}
\setlength{\partopsep}{0pt}
\item Phase08 run: \textbf{SmolVLA policy} rollouts.
\item Phase09 run: \textbf{oracle-action replay} baseline.
\item Metric: world-model latent distance $d\_goal\_l2\_wm\_int$.
\item Binning: env-action bins 0:5, 5:10, 10:15, 15:20, 20:25.
\item For each task, rows are shown as \textbf{VLA} and \textbf{$\Delta_{\mathrm{oracle}}$} (oracle minus VLA).
\item Task names have the \texttt{-v3} suffix removed for readability.
\end{itemize}

\begin{center}
\resizebox{\textwidth}{!}{%
\begin{tabular}{llrrrr|r|}
\hline
\textbf{Task} & \textbf{Run} & \textbf{0:5} & \textbf{5:10} & \textbf{10:15} & \textbf{15:20} & \textbf{20:25} \\
\hline
button-press-topdown & VLA & 509.30 & 480.13 & 418.32 & 350.63 & 331.17 \\
button-press-topdown & $\Delta_{\mathrm{oracle}}$ & +0.00 & -1.63 & -4.37 & -12.73 & -11.90 \\
\hline
door-open & VLA & 494.55 & 441.95 & 379.55 & 333.90 & 326.77 \\
door-open & $\Delta_{\mathrm{oracle}}$ & +0.03 & -0.42 & -4.97 & -15.70 & -28.22 \\
\hline
drawer-close & VLA & 434.17 & 435.65 & 437.80 & 440.60 & 441.52 \\
drawer-close & $\Delta_{\mathrm{oracle}}$ & -0.10 & -0.17 & -0.35 & -0.22 & +3.10 \\
\hline
drawer-open & VLA & 407.30 & 388.75 & 368.52 & 351.77 & 343.32 \\
drawer-open & $\Delta_{\mathrm{oracle}}$ & -2.38 & -6.43 & -12.52 & -19.52 & -26.42 \\
\hline
peg-insert-side & VLA & 454.18 & 434.30 & 415.48 & 405.15 & 404.77 \\
peg-insert-side & $\Delta_{\mathrm{oracle}}$ & -2.45 & -6.90 & -18.35 & -37.37 & -45.88 \\
\hline
pick-place & VLA & 381.93 & 378.55 & 370.47 & 353.80 & 339.18 \\
pick-place & $\Delta_{\mathrm{oracle}}$ & -11.68 & -23.60 & -27.32 & -26.37 & -18.92 \\
\hline
push & VLA & 341.73 & 337.92 & 338.53 & 333.77 & 332.22 \\
push & $\Delta_{\mathrm{oracle}}$ & -3.80 & -5.37 & -2.85 & -1.40 & -6.60 \\
\hline
reach & VLA & 418.10 & 382.05 & 353.47 & 337.63 & 341.57 \\
reach & $\Delta_{\mathrm{oracle}}$ & -0.68 & -4.50 & -12.70 & -13.33 & -9.77 \\
\hline
window-close & VLA & 572.08 & 544.30 & 502.33 & 435.35 & 346.00 \\
window-close & $\Delta_{\mathrm{oracle}}$ & +0.15 & +0.18 & +0.75 & +1.12 & -1.90 \\
\hline
window-open & VLA & 497.23 & 470.88 & 422.18 & 390.03 & 382.68 \\
window-open & $\Delta_{\mathrm{oracle}}$ & +0.05 & -2.33 & -8.00 & -18.13 & -16.60 \\
\hline
All tasks & VLA & 451.06 & 429.45 & 400.67 & 373.26 & 358.92 \\
All tasks & $\Delta_{\mathrm{oracle}}$ & -2.09 & -5.12 & -9.07 & -14.37 & -16.31 \\
\hline
\hline
\end{tabular}
}



% ensure this is a clean table and easily readable
% 3rd row for policy distance

\end{center}
\par\vspace{0.3em}
\noindent\small\textbf{How to read this table.} 9/10 tasks are closer to oracle rollout than the policy rollout for 25th frame, which is what we want. 
- Only \textit{drawer-close} isn't.
Improvement from the 7/10 on 50 frames ahead so this seems like a good action horizon.

Interestingly reducing the action horizon any further wouldn't lead to any improvements, as then "window-close" becomes a negative sample.

Cell values are cross-episode mean latent distances per bin. Each task has a VLA row and a \(\Delta_{\mathrm{oracle}}\) row, where \(\Delta_{\mathrm{oracle}} = \mathrm{oracle} - \mathrm{VLA}\). A positive value means oracle is farther from the world-model goal than VLA, and a negative value means oracle is closer. In this representation, larger raw values mean larger latent distance, so values closer to zero are preferable when the distance is interpreted as reconstruction error.


% SmolVLA uses MEAN/STD normalization for state and action; visual features use IDENTITY at the policy level.

% Dataset actions for MetaWorld MT50 are not pre‑clipped; they can exceed 
% [−1,1] in the stored data.

% Clipping to MetaWorld’s 
% [−1,1] action space happens only in the env_postprocessor at rollout time.




% Evidence:

% Legacy Phase11: max_steps=120.
% train_u1_to_u4: ~58.5 min for 4 updates = 14.6 min/update.
% long_to25: ~295.7 min for 20 updates = 14.8 min/update.
% Current Phase111 official:
% pilot 0→10: ~676 min for 10 updates = 67.6 min/update.
% resume 10→20: ~633 min for 10 updates = 63.3 min/update.
% Slowdown: ~4.3x to 4.6x vs legacy.
% CS238 report:

% Says Model 1: 20 updates took 12.26h, avg 0.613h/update = 36.8 min/update.
% Our current 63–68 min/update is ~1.7–1.8x slower than that.
% But their run = LIBERO, not our MetaWorld official LeRobot path. Compare useful, not exact.






\begin{table}[h]
\centering
\begin{tabular}{l l r r}
\hline
Run & Task & Episodes & Success (\%) \\
\hline
Non-fixed reset & pick-place-wall-v3 & 20 & 55.0 \\
Non-fixed reset & shelf-place-v3 & 20 & 30.0 \\
Non-fixed reset & stick-pull-v3 & 20 & 0.0 \\
Non-fixed reset & stick-push-v3 & 20 & 75.0 \\
\hline
Fixed reset & pick-place-wall-v3 & 20 & 0.0 \\
Fixed reset & shelf-place-v3 & 20 & 0.0 \\
Fixed reset & stick-pull-v3 & 20 & 0.0 \\
Fixed reset & stick-push-v3 & 20 & 30.0 \\
\hline
\end{tabular}
\caption{Phase072 hard-task evaluation, seed 1010, 20 episodes per task.}
\end{table}





% Check training data, ensure we use same tanh? / normalization method
\paragraph{SmolVLA Meta-World checkpoint contract.}
The public checkpoint \texttt{jadechoghari/smolvla\_metaworld} is a LeRobot SmolVLA policy fine-tuned from \texttt{lerobot/smolvla\_base} on the \texttt{lerobot/metaworld\_mt50} dataset, under Apache-2.0 license.\footnote{\url{https://huggingface.co/jadechoghari/smolvla_metaworld}}  Its public training configuration specifies Meta-World MT50 evaluation with RGB observations \texttt{observation.image}, proprioception \texttt{observation.state}, 4-D actions, seed \texttt{1000}, \texttt{chunk\_size=50}, \texttt{n\_action\_steps=1}, and synchronous eval environments (\texttt{use\_async\_envs=false}).  The normalization contract is \texttt{VISUAL: IDENTITY}, \texttt{STATE: MEAN\_STD}, and \texttt{ACTION: MEAN\_STD}; there is no config-level tanh action squashing contract in the official checkpoint metadata.\footnote{\url{https://huggingface.co/jadechoghari/smolvla_metaworld/raw/main/train_config.json}}  Therefore, for reproduction, the policy output path should preserve LeRobot's standard action unnormalization rather than assuming a tanh-squashed Gaussian action head.

% VLA 0 shot - MT10 % success rate
% check smolvla paper and reproduce results
\begin{table}[h]
  \centering
  \begin{tabular}{lrrrrr}
    \toprule
    Method & Easy & Medium & Hard & Very Hard & Avg. \\
    \midrule
    SmolVLA (0.45B), zero-shot Meta-World & 82.5 & 41.8 & 45.0 & 60.0 & 57.3 \\
    \bottomrule
  \end{tabular}
  \caption{Reported SmolVLA zero-shot Meta-World success rate by difficulty bucket.}
  \label{tab:smolvla-metaworld-reported-zero-shot}
\end{table}

% push-v3 15%

\paragraph{Meta-World evaluation ledger.}
Table~\ref{tab:mt50-eval-ledger} records the evaluation work used to debug the SmolVLA Meta-World reproduction.  The canonical path is now the official LeRobot evaluator with \texttt{jadechoghari/smolvla\_metaworld}; older custom runs are retained only as diagnostics because their wrappers and action handling differ from the current evaluation stack.

\begin{table}[h]
  \centering
  \small
  \begin{tabular}{p{0.22\linewidth}p{0.17\linewidth}p{0.17\linewidth}p{0.32\linewidth}}
    \toprule
    Run & Method / params & Result & Why kept \\
    \midrule
    \texttt{MT50\_Phase071\_official\_lerobot\_1task\_1ep} & Official LeRobot, assembly-v3, 1 episode, seed 1000, video on & 1/1 success & Smoke test: verified checkpoint load, environment launch, video writing, and success extraction. \\
    \texttt{MT50\_Phase072\_official\_lerobot\_10ep} & Official LeRobot, all 50 tasks, 10 episodes each, seed 1000, 3 GPU shards, default random resets & Difficulty-macro 45.0\%; very hard 17/50 & Main reproduction baseline.  Lower than reported SmolVLA macro 57.3\%, mainly very-hard tasks. \\
    \texttt{MT50\_Phase072\_official\_lerobot\_push\_10ep} & Official LeRobot, push-v3 only, 10 episodes, seed 1000, video off & 1/10 success & Checked no-video path and isolated slow single-task rerun. \\
    \texttt{MT50\_Phase073\_fixed\_veryhard\_10ep\_video} & Official LeRobot, five very-hard tasks, 10 episodes each, seed 1000, \texttt{MT50\_METAWORLD\_FREEZE\_RAND\_VEC=true}, video on & 8/50 success; disassemble-v3 8/10, others 0/10 & Tested fixed object/goal reset hypothesis.  Fixed reset improves determinism, but did not recover brittle grasp tasks. \\
    \texttt{MT50\_Phase074\_action\_debug\_3task\_3ep\_video} & Official LeRobot, shelf-place-v3/stick-pull-v3/stick-push-v3, 3 episodes each, random reset, no action transform, video and per-step action logs & 5/9 success; shelf-place 2/3, stick-pull 0/3, stick-push 3/3 & Debug run, not benchmark.  Logged normalized, postprocessed, and sent actions to test gripper/action-contract hypotheses. \\
    \texttt{MT50\_Phase17\_official\_lerobot\_10ep\_2gpu\_full} & Official LeRobot, all 50 tasks, 10 episodes each, seed 1000, two GPU shards, very-hard tasks first, video off & In progress / awaiting final \texttt{eval\_info.json} & Overnight clean full MT50 rerun with less video overhead. \\
    \texttt{MT50\_Phase07\_500\_1ep} and \texttt{MT50\_Phase07\_500\_10ep} & Earlier custom SmolVLA parity evaluator, 500-step episodes, seed 1000 & Diagnostic only & Useful history, but not the final benchmark path because wrappers/action handling differ from official LeRobot. \\
    \texttt{phase11\_env\_on\_policy\_grpo} & Custom on-policy GRPO stack and environment wrapper & Diagnostic only & Out of sync with official SmolVLA eval: includes custom policy wrapper/action treatment, so not directly comparable to LeRobot MT50 scores. \\
    \bottomrule
  \end{tabular}
  \caption{Evaluation runs used during the SmolVLA Meta-World reproduction audit.}
  \label{tab:mt50-eval-ledger}
\end{table}

\paragraph{Very-hard task check.}
The main failure mode is concentrated in the very-hard bucket.  Summing the randomized Phase072 run and the fixed-reset Phase073 run gives: disassemble-v3 12/20 (60\%), pick-place-wall-v3 3/20 (15\%), shelf-place-v3 1/20 (5\%), stick-pull-v3 0/20 (0\%), and stick-push-v3 9/20 (45\%).  Total very-hard success is 25/100 (25\%).  Thus fixed resets alone are not a sufficient explanation for the reported-performance gap.

\paragraph{Reset determinism.}
The \texttt{MT50\_METAWORLD\_FREEZE\_RAND\_VEC=true} flag is our launcher-level switch for a runtime adapter, not a permanent edit to LeRobot.  It keeps Meta-World's \texttt{\_freeze\_rand\_vec} enabled after LeRobot task setup, so object/goal positions are reused instead of resampled at every reset.  With default LeRobot, \texttt{\_freeze\_rand\_vec=false}; a fixed seed then gives a deterministic reset sequence only if task order, number of resets, process layout, wrappers, \texttt{eval.batch\_size=1}, and \texttt{eval.use\_async\_envs=false} are unchanged.  For exact multi-GPU comparability, the safer pattern is one task per worker process, unique output directory per task, stable task order, fixed seed, synchronous envs, and explicit logging of seed/reset mode.  This is reproducible, but not automatically paper-correct: Phase073 showed fixed resets can hurt several brittle tasks.

\paragraph{Action contract and gripper clipping.}
The SmolVLA checkpoint uses LeRobot \texttt{MEAN\_STD} action unnormalization, so postprocessed actions outside \([-1,1]\) are expected before the environment boundary.  Meta-World internally clips the XYZ end-effector part, but the gripper dimension is not equivalently clipped on our observed path.  Because the Gym action space is \texttt{Box(-1,1,shape=(4,))}, gripper values outside \([-1,1]\) violate the environment action contract.  Phase074 found frequent XYZ out-of-range values and low but real gripper out-of-range rate (up to about 3.7\% of steps on the debug tasks).  We therefore added a targeted ablation: keep XYZ unchanged, clip only the final gripper action via \texttt{MT50\_LEROBOT\_ACTION\_TRANSFORM=clip\_gripper}, and compare against a clean no-transform rerun before treating clipping as a benchmark change.














% table / diagram for VLA WM
% WM, VLA, env
% normalization, clipping, y/n
t





\section{Submission of papers to NeurIPS 2024}


Please read the instructions below carefully and follow them faithfully.


\subsection{Style}


Papers to be submitted to NeurIPS 2024 must be prepared according to the
instructions presented here. Papers may only be up to {\bf nine} pages long,
including figures. Additional pages \emph{containing only acknowledgments and
references} are allowed. Papers that exceed the page limit will not be
reviewed, or in any other way considered for presentation at the conference.


The margins in 2024 are the same as those in previous years.


Authors are required to use the NeurIPS \LaTeX{} style files obtainable at the
NeurIPS website as indicated below. Please make sure you use the current files
and not previous versions. Tweaking the style files may be grounds for
rejection.


\subsection{Retrieval of style files}


The style files for NeurIPS and other conference information are available on
the website at
\begin{center}
  \url{http://www.neurips.cc/}
\end{center}
The file \verb+neurips_2024.pdf+ contains these instructions and illustrates the
various formatting requirements your NeurIPS paper must satisfy.


The only supported style file for NeurIPS 2024 is \verb+neurips_2024.sty+,
rewritten for \LaTeXe{}.  \textbf{Previous style files for \LaTeX{} 2.09,
  Microsoft Word, and RTF are no longer supported!}


The \LaTeX{} style file contains three optional arguments: \verb+final+, which
creates a camera-ready copy, \verb+preprint+, which creates a preprint for
submission to, e.g., arXiv, and \verb+nonatbib+, which will not load the
\verb+natbib+ package for you in case of package clash.


\paragraph{Preprint option}
If you wish to post a preprint of your work online, e.g., on arXiv, using the
NeurIPS style, please use the \verb+preprint+ option. This will create a
nonanonymized version of your work with the text ``Preprint. Work in progress.''
in the footer. This version may be distributed as you see fit, as long as you do not say which conference it was submitted to. Please \textbf{do
  not} use the \verb+final+ option, which should \textbf{only} be used for
papers accepted to NeurIPS.


At submission time, please omit the \verb+final+ and \verb+preprint+
options. This will anonymize your submission and add line numbers to aid
review. Please do \emph{not} refer to these line numbers in your paper as they
will be removed during generation of camera-ready copies.


The file \verb+neurips_2024.tex+ may be used as a ``shell'' for writing your
paper. All you have to do is replace the author, title, abstract, and text of
the paper with your own.


The formatting instructions contained in these style files are summarized in
Sections \ref{gen_inst}, \ref{headings}, and \ref{others} below.


\section{General formatting instructions}
\label{gen_inst}


The text must be confined within a rectangle 5.5~inches (33~picas) wide and
9~inches (54~picas) long. The left margin is 1.5~inch (9~picas).  Use 10~point
type with a vertical spacing (leading) of 11~points.  Times New Roman is the
preferred typeface throughout, and will be selected for you by default.
Paragraphs are separated by \nicefrac{1}{2}~line space (5.5 points), with no
indentation.


The paper title should be 17~point, initial caps/lower case, bold, centered
between two horizontal rules. The top rule should be 4~points thick and the
bottom rule should be 1~point thick. Allow \nicefrac{1}{4}~inch space above and
below the title to rules. All pages should start at 1~inch (6~picas) from the
top of the page.


For the final version, authors' names are set in boldface, and each name is
centered above the corresponding address. The lead author's name is to be listed
first (left-most), and the co-authors' names (if different address) are set to
follow. If there is only one co-author, list both author and co-author side by
side.


Please pay special attention to the instructions in Section \ref{others}
regarding figures, tables, acknowledgments, and references.


\section{Headings: first level}
\label{headings}


All headings should be lower case (except for first word and proper nouns),
flush left, and bold.


First-level headings should be in 12-point type.


\subsection{Headings: second level}


Second-level headings should be in 10-point type.


\subsubsection{Headings: third level}


Third-level headings should be in 10-point type.


\paragraph{Paragraphs}


There is also a \verb+\paragraph+ command available, which sets the heading in
bold, flush left, and inline with the text, with the heading followed by 1\,em
of space.


\section{Citations, figures, tables, references}
\label{others}


These instructions apply to everyone.


\subsection{Citations within the text}


The \verb+natbib+ package will be loaded for you by default.  Citations may be
author/year or numeric, as long as you maintain internal consistency.  As to the
format of the references themselves, any style is acceptable as long as it is
used consistently.


The documentation for \verb+natbib+ may be found at
\begin{center}
  \url{http://mirrors.ctan.org/macros/latex/contrib/natbib/natnotes.pdf}
\end{center}
Of note is the command \verb+\citet+, which produces citations appropriate for
use in inline text.  For example,
\begin{verbatim}
   \citet{hasselmo} investigated\dots
\end{verbatim}
produces
\begin{quote}
  Hasselmo, et al.\ (1995) investigated\dots
\end{quote}


If you wish to load the \verb+natbib+ package with options, you may add the
following before loading the \verb+neurips_2024+ package:
\begin{verbatim}
   \PassOptionsToPackage{options}{natbib}
\end{verbatim}


If \verb+natbib+ clashes with another package you load, you can add the optional
argument \verb+nonatbib+ when loading the style file:
\begin{verbatim}
   \usepackage[nonatbib]{neurips_2024}
\end{verbatim}


As submission is double blind, refer to your own published work in the third
person. That is, use ``In the previous work of Jones et al.\ [4],'' not ``In our
previous work [4].'' If you cite your other papers that are not widely available
(e.g., a journal paper under review), use anonymous author names in the
citation, e.g., an author of the form ``A.\ Anonymous'' and include a copy of the anonymized paper in the supplementary material.


\subsection{Footnotes}


Footnotes should be used sparingly.  If you do require a footnote, indicate
footnotes with a number\footnote{Sample of the first footnote.} in the
text. Place the footnotes at the bottom of the page on which they appear.
Precede the footnote with a horizontal rule of 2~inches (12~picas).


Note that footnotes are properly typeset \emph{after} punctuation
marks.\footnote{As in this example.}


\subsection{Figures}


\begin{figure}
  \centering
  \fbox{\rule[-.5cm]{0cm}{4cm} \rule[-.5cm]{4cm}{0cm}}
  \caption{Sample figure caption.}
\end{figure}


All artwork must be neat, clean, and legible. Lines should be dark enough for
purposes of reproduction. The figure number and caption always appear after the
figure. Place one line space before the figure caption and one line space after
the figure. The figure caption should be lower case (except for first word and
proper nouns); figures are numbered consecutively.


You may use color figures.  However, it is best for the figure captions and the
paper body to be legible if the paper is printed in either black/white or in
color.


\subsection{Tables}


All tables must be centered, neat, clean and legible.  The table number and
title always appear before the table.  See Table~\ref{sample-table}.


Place one line space before the table title, one line space after the
table title, and one line space after the table. The table title must
be lower case (except for first word and proper nouns); tables are
numbered consecutively.


Note that publication-quality tables \emph{do not contain vertical rules.} We
strongly suggest the use of the \verb+booktabs+ package, which allows for
typesetting high-quality, professional tables:
\begin{center}
  \url{https://www.ctan.org/pkg/booktabs}
\end{center}
This package was used to typeset Table~\ref{sample-table}.


\begin{table}
  \caption{Sample table title}
  \label{sample-table}
  \centering
  \begin{tabular}{lll}
    \toprule
    \multicolumn{2}{c}{Part}                   \\
    \cmidrule(r){1-2}
    Name     & Description     & Size ($\mu$m) \\
    \midrule
    Dendrite & Input terminal  & $\sim$100     \\
    Axon     & Output terminal & $\sim$10      \\
    Soma     & Cell body       & up to $10^6$  \\
    \bottomrule
  \end{tabular}
\end{table}

\subsection{Math}
Note that display math in bare TeX commands will not create correct line numbers for submission. Please use LaTeX (or AMSTeX) commands for unnumbered display math. (You really shouldn't be using \$\$ anyway; see \url{https://tex.stackexchange.com/questions/503/why-is-preferable-to} and \url{https://tex.stackexchange.com/questions/40492/what-are-the-differences-between-align-equation-and-displaymath} for more information.)

\subsection{Final instructions}

Do not change any aspects of the formatting parameters in the style files.  In
particular, do not modify the width or length of the rectangle the text should
fit into, and do not change font sizes (except perhaps in the
\textbf{References} section; see below). Please note that pages should be
numbered.


\section{Preparing PDF files}


Please prepare submission files with paper size ``US Letter,'' and not, for
example, ``A4.''


Fonts were the main cause of problems in the past years. Your PDF file must only
contain Type 1 or Embedded TrueType fonts. Here are a few instructions to
achieve this.


\begin{itemize}


\item You should directly generate PDF files using \verb+pdflatex+.


\item You can check which fonts a PDF files uses.  In Acrobat Reader, select the
  menu Files$>$Document Properties$>$Fonts and select Show All Fonts. You can
  also use the program \verb+pdffonts+ which comes with \verb+xpdf+ and is
  available out-of-the-box on most Linux machines.


\item \verb+xfig+ "patterned" shapes are implemented with bitmap fonts.  Use
  "solid" shapes instead.


\item The \verb+\bbold+ package almost always uses bitmap fonts.  You should use
  the equivalent AMS Fonts:
\begin{verbatim}
   \usepackage{amsfonts}
\end{verbatim}
followed by, e.g., \verb+\mathbb{R}+, \verb+\mathbb{N}+, or \verb+\mathbb{C}+
for $\mathbb{R}$, $\mathbb{N}$ or $\mathbb{C}$.  You can also use the following
workaround for reals, natural and complex:
\begin{verbatim}
   \newcommand{\RR}{I\!\!R} %real numbers
   \newcommand{\Nat}{I\!\!N} %natural numbers
   \newcommand{\CC}{I\!\!\!\!C} %complex numbers
\end{verbatim}
Note that \verb+amsfonts+ is automatically loaded by the \verb+amssymb+ package.


\end{itemize}


If your file contains type 3 fonts or non embedded TrueType fonts, we will ask
you to fix it.


\subsection{Margins in \LaTeX{}}


Most of the margin problems come from figures positioned by hand using
\verb+\special+ or other commands. We suggest using the command
\verb+\includegraphics+ from the \verb+graphicx+ package. Always specify the
figure width as a multiple of the line width as in the example below:
\begin{verbatim}
   \usepackage[pdftex]{graphicx} ...
   \includegraphics[width=0.8\linewidth]{myfile.pdf}
\end{verbatim}
See Section 4.4 in the graphics bundle documentation
(\url{http://mirrors.ctan.org/macros/latex/required/graphics/grfguide.pdf})


A number of width problems arise when \LaTeX{} cannot properly hyphenate a
line. Please give LaTeX hyphenation hints using the \verb+\-+ command when
necessary.

\begin{ack}
Use unnumbered first level headings for the acknowledgments. All acknowledgments
go at the end of the paper before the list of references. Moreover, you are required to declare
funding (financial activities supporting the submitted work) and competing interests (related financial activities outside the submitted work).
More information about this disclosure can be found at: \url{https://neurips.cc/Conferences/2024/PaperInformation/FundingDisclosure}.


Do {\bf not} include this section in the anonymized submission, only in the final paper. You can use the \texttt{ack} environment provided in the style file to automatically hide this section in the anonymized submission.
\end{ack}

\section*{References}


References follow the acknowledgments in the camera-ready paper. Use unnumbered first-level heading for
the references. Any choice of citation style is acceptable as long as you are
consistent. It is permissible to reduce the font size to \verb+small+ (9 point)
when listing the references.
Note that the Reference section does not count towards the page limit.
\medskip


{
\small


[1] Alexander, J.A.\ \& Mozer, M.C.\ (1995) Template-based algorithms for
connectionist rule extraction. In G.\ Tesauro, D.S.\ Touretzky and T.K.\ Leen
(eds.), {\it Advances in Neural Information Processing Systems 7},
pp.\ 609--616. Cambridge, MA: MIT Press.


[2] Bower, J.M.\ \& Beeman, D.\ (1995) {\it The Book of GENESIS: Exploring
  Realistic Neural Models with the GEneral NEural SImulation System.}  New York:
TELOS/Springer--Verlag.


[3] Hasselmo, M.E., Schnell, E.\ \& Barkai, E.\ (1995) Dynamics of learning and
recall at excitatory recurrent synapses and cholinergic modulation in rat
hippocampal region CA3. {\it Journal of Neuroscience} {\bf 15}(7):5249-5262.
}


%%%%%%%%%%%%%%%%%%%%%%%%%%%%%%%%%%%%%%%%%%%%%%%%%%%%%%%%%%%%

\appendix

\section{Appendix / supplemental material}


Optionally include supplemental material (complete proofs, additional experiments and plots) in appendix.
All such materials \textbf{SHOULD be included in the main submission.}

%%%%%%%%%%%%%%%%%%%%%%%%%%%%%%%%%%%%%%%%%%%%%%%%%%%%%%%%%%%%

\newpage
\section*{NeurIPS Paper Checklist}

%%% BEGIN INSTRUCTIONS %%%
The checklist is designed to encourage best practices for responsible machine learning research, addressing issues of reproducibility, transparency, research ethics, and societal impact. Do not remove the checklist: {\bf The papers not including the checklist will be desk rejected.} The checklist should follow the references and follow the (optional) supplemental material.  The checklist does NOT count towards the page
limit. 

Please read the checklist guidelines carefully for information on how to answer these questions. For each question in the checklist:
\begin{itemize}
    \item You should answer \answerYes{}, \answerNo{}, or \answerNA{}.
    \item \answerNA{} means either that the question is Not Applicable for that particular paper or the relevant information is Not Available.
    \item Please provide a short (1–2 sentence) justification right after your answer (even for NA). 
   % \item {\bf The papers not including the checklist will be desk rejected.}
\end{itemize}

{\bf The checklist answers are an integral part of your paper submission.} They are visible to the reviewers, area chairs, senior area chairs, and ethics reviewers. You will be asked to also include it (after eventual revisions) with the final version of your paper, and its final version will be published with the paper.

The reviewers of your paper will be asked to use the checklist as one of the factors in their evaluation. While "\answerYes{}" is generally preferable to "\answerNo{}", it is perfectly acceptable to answer "\answerNo{}" provided a proper justification is given (e.g., "error bars are not reported because it would be too computationally expensive" or "we were unable to find the license for the dataset we used"). In general, answering "\answerNo{}" or "\answerNA{}" is not grounds for rejection. While the questions are phrased in a binary way, we acknowledge that the true answer is often more nuanced, so please just use your best judgment and write a justification to elaborate. All supporting evidence can appear either in the main paper or the supplemental material, provided in appendix. If you answer \answerYes{} to a question, in the justification please point to the section(s) where related material for the question can be found.

IMPORTANT, please:
\begin{itemize}
    \item {\bf Delete this instruction block, but keep the section heading ``NeurIPS paper checklist"},
    \item  {\bf Keep the checklist subsection headings, questions/answers and guidelines below.}
    \item {\bf Do not modify the questions and only use the provided macros for your answers}.
\end{itemize} 
 

%%% END INSTRUCTIONS %%%


\begin{enumerate}

\item {\bf Claims}
    \item[] Question: Do the main claims made in the abstract and introduction accurately reflect the paper's contributions and scope?
    \item[] Answer: \answerTODO{} % Replace by \answerYes{}, \answerNo{}, or \answerNA{}.
    \item[] Justification: \justificationTODO{}
    \item[] Guidelines:
    \begin{itemize}
        \item The answer NA means that the abstract and introduction do not include the claims made in the paper.
        \item The abstract and/or introduction should clearly state the claims made, including the contributions made in the paper and important assumptions and limitations. A No or NA answer to this question will not be perceived well by the reviewers. 
        \item The claims made should match theoretical and experimental results, and reflect how much the results can be expected to generalize to other settings. 
        \item It is fine to include aspirational goals as motivation as long as it is clear that these goals are not attained by the paper. 
    \end{itemize}

\item {\bf Limitations}
    \item[] Question: Does the paper discuss the limitations of the work performed by the authors?
    \item[] Answer: \answerTODO{} % Replace by \answerYes{}, \answerNo{}, or \answerNA{}.
    \item[] Justification: \justificationTODO{}
    \item[] Guidelines:
    \begin{itemize}
        \item The answer NA means that the paper has no limitation while the answer No means that the paper has limitations, but those are not discussed in the paper. 
        \item The authors are encouraged to create a separate "Limitations" section in their paper.
        \item The paper should point out any strong assumptions and how robust the results are to violations of these assumptions (e.g., independence assumptions, noiseless settings, model well-specification, asymptotic approximations only holding locally). The authors should reflect on how these assumptions might be violated in practice and what the implications would be.
        \item The authors should reflect on the scope of the claims made, e.g., if the approach was only tested on a few datasets or with a few runs. In general, empirical results often depend on implicit assumptions, which should be articulated.
        \item The authors should reflect on the factors that influence the performance of the approach. For example, a facial recognition algorithm may perform poorly when image resolution is low or images are taken in low lighting. Or a speech-to-text system might not be used reliably to provide closed captions for online lectures because it fails to handle technical jargon.
        \item The authors should discuss the computational efficiency of the proposed algorithms and how they scale with dataset size.
        \item If applicable, the authors should discuss possible limitations of their approach to address problems of privacy and fairness.
        \item While the authors might fear that complete honesty about limitations might be used by reviewers as grounds for rejection, a worse outcome might be that reviewers discover limitations that aren't acknowledged in the paper. The authors should use their best judgment and recognize that individual actions in favor of transparency play an important role in developing norms that preserve the integrity of the community. Reviewers will be specifically instructed to not penalize honesty concerning limitations.
    \end{itemize}

\item {\bf Theory Assumptions and Proofs}
    \item[] Question: For each theoretical result, does the paper provide the full set of assumptions and a complete (and correct) proof?
    \item[] Answer: \answerTODO{} % Replace by \answerYes{}, \answerNo{}, or \answerNA{}.
    \item[] Justification: \justificationTODO{}
    \item[] Guidelines:
    \begin{itemize}
        \item The answer NA means that the paper does not include theoretical results. 
        \item All the theorems, formulas, and proofs in the paper should be numbered and cross-referenced.
        \item All assumptions should be clearly stated or referenced in the statement of any theorems.
        \item The proofs can either appear in the main paper or the supplemental material, but if they appear in the supplemental material, the authors are encouraged to provide a short proof sketch to provide intuition. 
        \item Inversely, any informal proof provided in the core of the paper should be complemented by formal proofs provided in appendix or supplemental material.
        \item Theorems and Lemmas that the proof relies upon should be properly referenced. 
    \end{itemize}

    \item {\bf Experimental Result Reproducibility}
    \item[] Question: Does the paper fully disclose all the information needed to reproduce the main experimental results of the paper to the extent that it affects the main claims and/or conclusions of the paper (regardless of whether the code and data are provided or not)?
    \item[] Answer: \answerTODO{} % Replace by \answerYes{}, \answerNo{}, or \answerNA{}.
    \item[] Justification: \justificationTODO{}
    \item[] Guidelines:
    \begin{itemize}
        \item The answer NA means that the paper does not include experiments.
        \item If the paper includes experiments, a No answer to this question will not be perceived well by the reviewers: Making the paper reproducible is important, regardless of whether the code and data are provided or not.
        \item If the contribution is a dataset and/or model, the authors should describe the steps taken to make their results reproducible or verifiable. 
        \item Depending on the contribution, reproducibility can be accomplished in various ways. For example, if the contribution is a novel architecture, describing the architecture fully might suffice, or if the contribution is a specific model and empirical evaluation, it may be necessary to either make it possible for others to replicate the model with the same dataset, or provide access to the model. In general. releasing code and data is often one good way to accomplish this, but reproducibility can also be provided via detailed instructions for how to replicate the results, access to a hosted model (e.g., in the case of a large language model), releasing of a model checkpoint, or other means that are appropriate to the research performed.
        \item While NeurIPS does not require releasing code, the conference does require all submissions to provide some reasonable avenue for reproducibility, which may depend on the nature of the contribution. For example
        \begin{enumerate}
            \item If the contribution is primarily a new algorithm, the paper should make it clear how to reproduce that algorithm.
            \item If the contribution is primarily a new model architecture, the paper should describe the architecture clearly and fully.
            \item If the contribution is a new model (e.g., a large language model), then there should either be a way to access this model for reproducing the results or a way to reproduce the model (e.g., with an open-source dataset or instructions for how to construct the dataset).
            \item We recognize that reproducibility may be tricky in some cases, in which case authors are welcome to describe the particular way they provide for reproducibility. In the case of closed-source models, it may be that access to the model is limited in some way (e.g., to registered users), but it should be possible for other researchers to have some path to reproducing or verifying the results.
        \end{enumerate}
    \end{itemize}


\item {\bf Open access to data and code}
    \item[] Question: Does the paper provide open access to the data and code, with sufficient instructions to faithfully reproduce the main experimental results, as described in supplemental material?
    \item[] Answer: \answerTODO{} % Replace by \answerYes{}, \answerNo{}, or \answerNA{}.
    \item[] Justification: \justificationTODO{}
    \item[] Guidelines:
    \begin{itemize}
        \item The answer NA means that paper does not include experiments requiring code.
        \item Please see the NeurIPS code and data submission guidelines (\url{https://nips.cc/public/guides/CodeSubmissionPolicy}) for more details.
        \item While we encourage the release of code and data, we understand that this might not be possible, so “No” is an acceptable answer. Papers cannot be rejected simply for not including code, unless this is central to the contribution (e.g., for a new open-source benchmark).
        \item The instructions should contain the exact command and environment needed to run to reproduce the results. See the NeurIPS code and data submission guidelines (\url{https://nips.cc/public/guides/CodeSubmissionPolicy}) for more details.
        \item The authors should provide instructions on data access and preparation, including how to access the raw data, preprocessed data, intermediate data, and generated data, etc.
        \item The authors should provide scripts to reproduce all experimental results for the new proposed method and baselines. If only a subset of experiments are reproducible, they should state which ones are omitted from the script and why.
        \item At submission time, to preserve anonymity, the authors should release anonymized versions (if applicable).
        \item Providing as much information as possible in supplemental material (appended to the paper) is recommended, but including URLs to data and code is permitted.
    \end{itemize}


\item {\bf Experimental Setting/Details}
    \item[] Question: Does the paper specify all the training and test details (e.g., data splits, hyperparameters, how they were chosen, type of optimizer, etc.) necessary to understand the results?
    \item[] Answer: \answerTODO{} % Replace by \answerYes{}, \answerNo{}, or \answerNA{}.
    \item[] Justification: \justificationTODO{}
    \item[] Guidelines:
    \begin{itemize}
        \item The answer NA means that the paper does not include experiments.
        \item The experimental setting should be presented in the core of the paper to a level of detail that is necessary to appreciate the results and make sense of them.
        \item The full details can be provided either with the code, in appendix, or as supplemental material.
    \end{itemize}

\item {\bf Experiment Statistical Significance}
    \item[] Question: Does the paper report error bars suitably and correctly defined or other appropriate information about the statistical significance of the experiments?
    \item[] Answer: \answerTODO{} % Replace by \answerYes{}, \answerNo{}, or \answerNA{}.
    \item[] Justification: \justificationTODO{}
    \item[] Guidelines:
    \begin{itemize}
        \item The answer NA means that the paper does not include experiments.
        \item The authors should answer "Yes" if the results are accompanied by error bars, confidence intervals, or statistical significance tests, at least for the experiments that support the main claims of the paper.
        \item The factors of variability that the error bars are capturing should be clearly stated (for example, train/test split, initialization, random drawing of some parameter, or overall run with given experimental conditions).
        \item The method for calculating the error bars should be explained (closed form formula, call to a library function, bootstrap, etc.)
        \item The assumptions made should be given (e.g., Normally distributed errors).
        \item It should be clear whether the error bar is the standard deviation or the standard error of the mean.
        \item It is OK to report 1-sigma error bars, but one should state it. The authors should preferably report a 2-sigma error bar than state that they have a 96\% CI, if the hypothesis of Normality of errors is not verified.
        \item For asymmetric distributions, the authors should be careful not to show in tables or figures symmetric error bars that would yield results that are out of range (e.g. negative error rates).
        \item If error bars are reported in tables or plots, The authors should explain in the text how they were calculated and reference the corresponding figures or tables in the text.
    \end{itemize}

\item {\bf Experiments Compute Resources}
    \item[] Question: For each experiment, does the paper provide sufficient information on the computer resources (type of compute workers, memory, time of execution) needed to reproduce the experiments?
    \item[] Answer: \answerTODO{} % Replace by \answerYes{}, \answerNo{}, or \answerNA{}.
    \item[] Justification: \justificationTODO{}
    \item[] Guidelines:
    \begin{itemize}
        \item The answer NA means that the paper does not include experiments.
        \item The paper should indicate the type of compute workers CPU or GPU, internal cluster, or cloud provider, including relevant memory and storage.
        \item The paper should provide the amount of compute required for each of the individual experimental runs as well as estimate the total compute. 
        \item The paper should disclose whether the full research project required more compute than the experiments reported in the paper (e.g., preliminary or failed experiments that didn't make it into the paper). 
    \end{itemize}
    
\item {\bf Code Of Ethics}
    \item[] Question: Does the research conducted in the paper conform, in every respect, with the NeurIPS Code of Ethics \url{https://neurips.cc/public/EthicsGuidelines}?
    \item[] Answer: \answerTODO{} % Replace by \answerYes{}, \answerNo{}, or \answerNA{}.
    \item[] Justification: \justificationTODO{}
    \item[] Guidelines:
    \begin{itemize}
        \item The answer NA means that the authors have not reviewed the NeurIPS Code of Ethics.
        \item If the authors answer No, they should explain the special circumstances that require a deviation from the Code of Ethics.
        \item The authors should make sure to preserve anonymity (e.g., if there is a special consideration due to laws or regulations in their jurisdiction).
    \end{itemize}


\item {\bf Broader Impacts}
    \item[] Question: Does the paper discuss both potential positive societal impacts and negative societal impacts of the work performed?
    \item[] Answer: \answerTODO{} % Replace by \answerYes{}, \answerNo{}, or \answerNA{}.
    \item[] Justification: \justificationTODO{}
    \item[] Guidelines:
    \begin{itemize}
        \item The answer NA means that there is no societal impact of the work performed.
        \item If the authors answer NA or No, they should explain why their work has no societal impact or why the paper does not address societal impact.
        \item Examples of negative societal impacts include potential malicious or unintended uses (e.g., disinformation, generating fake profiles, surveillance), fairness considerations (e.g., deployment of technologies that could make decisions that unfairly impact specific groups), privacy considerations, and security considerations.
        \item The conference expects that many papers will be foundational research and not tied to particular applications, let alone deployments. However, if there is a direct path to any negative applications, the authors should point it out. For example, it is legitimate to point out that an improvement in the quality of generative models could be used to generate deepfakes for disinformation. On the other hand, it is not needed to point out that a generic algorithm for optimizing neural networks could enable people to train models that generate Deepfakes faster.
        \item The authors should consider possible harms that could arise when the technology is being used as intended and functioning correctly, harms that could arise when the technology is being used as intended but gives incorrect results, and harms following from (intentional or unintentional) misuse of the technology.
        \item If there are negative societal impacts, the authors could also discuss possible mitigation strategies (e.g., gated release of models, providing defenses in addition to attacks, mechanisms for monitoring misuse, mechanisms to monitor how a system learns from feedback over time, improving the efficiency and accessibility of ML).
    \end{itemize}
    
\item {\bf Safeguards}
    \item[] Question: Does the paper describe safeguards that have been put in place for responsible release of data or models that have a high risk for misuse (e.g., pretrained language models, image generators, or scraped datasets)?
    \item[] Answer: \answerTODO{} % Replace by \answerYes{}, \answerNo{}, or \answerNA{}.
    \item[] Justification: \justificationTODO{}
    \item[] Guidelines:
    \begin{itemize}
        \item The answer NA means that the paper poses no such risks.
        \item Released models that have a high risk for misuse or dual-use should be released with necessary safeguards to allow for controlled use of the model, for example by requiring that users adhere to usage guidelines or restrictions to access the model or implementing safety filters. 
        \item Datasets that have been scraped from the Internet could pose safety risks. The authors should describe how they avoided releasing unsafe images.
        \item We recognize that providing effective safeguards is challenging, and many papers do not require this, but we encourage authors to take this into account and make a best faith effort.
    \end{itemize}

\item {\bf Licenses for existing assets}
    \item[] Question: Are the creators or original owners of assets (e.g., code, data, models), used in the paper, properly credited and are the license and terms of use explicitly mentioned and properly respected?
    \item[] Answer: \answerTODO{} % Replace by \answerYes{}, \answerNo{}, or \answerNA{}.
    \item[] Justification: \justificationTODO{}
    \item[] Guidelines:
    \begin{itemize}
        \item The answer NA means that the paper does not use existing assets.
        \item The authors should cite the original paper that produced the code package or dataset.
        \item The authors should state which version of the asset is used and, if possible, include a URL.
        \item The name of the license (e.g., CC-BY 4.0) should be included for each asset.
        \item For scraped data from a particular source (e.g., website), the copyright and terms of service of that source should be provided.
        \item If assets are released, the license, copyright information, and terms of use in the package should be provided. For popular datasets, \url{paperswithcode.com/datasets} has curated licenses for some datasets. Their licensing guide can help determine the license of a dataset.
        \item For existing datasets that are re-packaged, both the original license and the license of the derived asset (if it has changed) should be provided.
        \item If this information is not available online, the authors are encouraged to reach out to the asset's creators.
    \end{itemize}

\item {\bf New Assets}
    \item[] Question: Are new assets introduced in the paper well documented and is the documentation provided alongside the assets?
    \item[] Answer: \answerTODO{} % Replace by \answerYes{}, \answerNo{}, or \answerNA{}.
    \item[] Justification: \justificationTODO{}
    \item[] Guidelines:
    \begin{itemize}
        \item The answer NA means that the paper does not release new assets.
        \item Researchers should communicate the details of the dataset/code/model as part of their submissions via structured templates. This includes details about training, license, limitations, etc. 
        \item The paper should discuss whether and how consent was obtained from people whose asset is used.
        \item At submission time, remember to anonymize your assets (if applicable). You can either create an anonymized URL or include an anonymized zip file.
    \end{itemize}

\item {\bf Crowdsourcing and Research with Human Subjects}
    \item[] Question: For crowdsourcing experiments and research with human subjects, does the paper include the full text of instructions given to participants and screenshots, if applicable, as well as details about compensation (if any)? 
    \item[] Answer: \answerTODO{} % Replace by \answerYes{}, \answerNo{}, or \answerNA{}.
    \item[] Justification: \justificationTODO{}
    \item[] Guidelines:
    \begin{itemize}
        \item The answer NA means that the paper does not involve crowdsourcing nor research with human subjects.
        \item Including this information in the supplemental material is fine, but if the main contribution of the paper involves human subjects, then as much detail as possible should be included in the main paper. 
        \item According to the NeurIPS Code of Ethics, workers involved in data collection, curation, or other labor should be paid at least the minimum wage in the country of the data collector. 
    \end{itemize}

\item {\bf Institutional Review Board (IRB) Approvals or Equivalent for Research with Human Subjects}
    \item[] Question: Does the paper describe potential risks incurred by study participants, whether such risks were disclosed to the subjects, and whether Institutional Review Board (IRB) approvals (or an equivalent approval/review based on the requirements of your country or institution) were obtained?
    \item[] Answer: \answerTODO{} % Replace by \answerYes{}, \answerNo{}, or \answerNA{}.
    \item[] Justification: \justificationTODO{}
    \item[] Guidelines:
    \begin{itemize}
        \item The answer NA means that the paper does not involve crowdsourcing nor research with human subjects.
        \item Depending on the country in which research is conducted, IRB approval (or equivalent) may be required for any human subjects research. If you obtained IRB approval, you should clearly state this in the paper. 
        \item We recognize that the procedures for this may vary significantly between institutions and locations, and we expect authors to adhere to the NeurIPS Code of Ethics and the guidelines for their institution. 
        \item For initial submissions, do not include any information that would break anonymity (if applicable), such as the institution conducting the review.
    \end{itemize}

\end{enumerate}


\end{document}
